% A.74 reserved by the root.
\subsection{The full congruence comparison on every odd-prime cusp disc}

\begin{statement}{A.74}{Proved}
{An explicit analytic point-to-cusp map with the exact separating Kummer level}
Let $p$ be an odd prime.  There is an explicit analytic isometry
and bijection
\[
 \Theta_p:p\Z_p\longrightarrow p\Z_p,\qquad
 \Theta_p(t)=t\sqrt{F(t)}=16\sqrt{q(t)},
\tag{A.74.1}
\]
where $q(t)=t^2F(t)/256$ is the formal Legendre Tate parameter,
and the square root is the one with constant term one.
For every nonzero $x\in p\Z_p$, the point section $s_x$ and
the cusp-zero section with tangent $\Theta_p(x)$ have
geometrically conjugate images in the full Legendre
congruence extension over $G_{\Q_p}$, with conjugator in $H$.

Put $e=v_p(x)>0$.  Their abelian Kummer images modulo $p^k$
agree for $0\le k<e$.  At level $p^e$ they differ, and the
specified finite etale cover
\[
                         z^{p^e}=1-t
\tag{A.74.2}
\]
lifts the cusp section but not $s_x$.
This cover does not factor through $H$.
Indeed the geometric pro-$p$ projection of $s_x$ lies in none
of the three cusp packets, already in its abelian quotient
modulo $p^e$.
The first separating level just asserted concerns these
abelian $p$-power Kummer images.

If $x$ is rational, the ordinary local tangent $\Theta_p(x)$
is transcendental over $\Q$.
All these assertions concern local sections over $\Q_p$.
\end{statement}
\begin{proof}
The formal inversion used in A.63 and A.67 gives
\[
 q(t)=\frac{t^2}{256}F(t),\qquad
 F(t)\in1+t\Z[1/2][[t]],\qquad
                        F(t)=1+t+O(t^2).
\tag{A.74.3}
\]
For the last coefficient, expand
\[
 \frac1{j(E_t)}=
  \frac{t^2}{256}\frac{(1-t)^2}{(t^2-t+1)^3}
                  =\frac{t^2}{256}(1+t+O(t^2)),
\]
and use $q=1/j+O(1/j^2)$.
The unique square root with constant term one is obtained
recursively by dividing each successive coefficient by two.
It consequently belongs to $1+t\Z[1/2][[t]]$ and begins
\[
                         \sqrt{F(t)}=1+t/2+O(t^2).
\tag{A.74.4}
\]
These series have integral coefficients in $\Z_p$ and
converge on $p\Z_p$.
In particular
$\Theta_p(t)\in t+t^2\Z_p[[t]]$.
For $x,y\in p\Z_p$, factoring every difference of powers
gives
\[
                \Theta_p(x)-\Theta_p(y)
                         \in(x-y)(1+p\Z_p).
\tag{A.74.5}
\]
This proves the isometry.  For any prescribed $C\in p\Z_p$,
the map $t\mapsto C-(\Theta_p(t)-t)$ is a strict contraction
of $p\Z_p$ into itself.  Its unique fixed point maps to $C$.
Thus (A.74.1) is a bijection and maps zero only to zero.

Let $C=\Theta_p(x)$ and $q_x=q(x)$.
They are actual ordinary local numbers satisfying
\[
                C^2/256=q_x,\qquad
                v_p(q_x)=2e,\qquad
                C/x\in1+p\Z_p.
\tag{A.74.6}
\]
The Legendre curve $E_x$ has unit $c_4$ and positive
discriminant valuation $2e$, so is multiplicative.
Its nodal slopes at zero are $\pm\sqrt{-1}$.
Its quadratic splitting character is therefore exactly
the character of $-1$, which is also the constant
splitting character of the cusp-zero representation.
After twisting both by this same character, the point
is the split Tate curve with parameter $q_x$ and the
cusp's integral Tate extension has class
$\kappa(C^2/256)=\kappa(q_x)$.
The formal unit in (A.74.3) has compatible roots of
constant term one on the formal cusp, as in A.63;
those roots are not evaluated at $x$ in this argument.
In particular no assertion that their $p$-power roots
converge at $x$ enters the coefficient-$p$ comparison.

Here is the exact intertwiner, including the arithmetic
action.  Choose compatible roots of unity and the usual
root-of-parameter Tate bases, paired by
$e_N(\zeta_N^a q^{b/N},\zeta_N^c q^{d/N})
=\zeta_N^{ad-bc}$ as in A.67.  Write
$\eta=\chi_{-1}:G_{\Q_p}\to\{1,-1\}$ for the common
splitting character.  In these bases the point and cusp
representations have matrices
\[
 \rho_i(g)=\eta(g)
    \begin{pmatrix}\chi_\ell(g)&c_i(g)\\0&1\end{pmatrix},
 \qquad c_2(g)-c_1(g)=(\chi_\ell(g)-1)h_\ell,
       \qquad h_\ell\in\Z_\ell.
\tag{A.74.6a}
\]
The integral zero-cochain $h_\ell$ exists because both
extension classes are $\kappa_\ell(q_x)$.
For $U_\ell=\left(\begin{smallmatrix}1&h_\ell\\0&1
\end{smallmatrix}\right)$ direct multiplication gives
$\rho_2=U_\ell^{-1}\rho_1U_\ell$.
The curve-to-Tate isomorphisms and these paired bases
preserve the Weil pairing; so does $U_\ell$, whose
determinant is one.  Undoing the same scalar twist
$\eta$ preserves this equality.  Thus at every odd
coefficient prime, including $p$, the resulting
conjugator belongs to
$\operatorname{SL}_2(\Z_\ell)=H_\ell$.

At two use the actual prime-to-$p$ formal comparison.
The ring and evaluation map are
\[
 R_p=\Z_p[[t]][1/(pt)],\qquad
 R_p\longrightarrow\Q_p,
                         \quad t\longmapsto x.
\tag{A.74.7}
\]
The map is defined because every power series converges
at $x$ and the denominators in an element of $R_p$
are bounded powers of $p$ and $t$.
The tame normalization, purity and finite-constant
descent proof in A.49 applies to this ring for every
odd $p$.  In its geometric pro-two quotient, every
cover is dominated after a finite constant extension
by $t\mapsto t^{2^a}$.
Thus the geometric group is $\Z_2(1)$ with the
cyclotomic arithmetic action; comparison with
$\Q_p((t))$ identifies it with the pro-two
cusp-zero decomposition group.
Pulling back $z^{2^a}=t$ along (A.74.7) shows that
the evaluation section has precisely the tangent
class $\kappa_2(x)$.
This gives an actual geometric pro-two conjugacy of
$s_x$ with that cusp section in the pro-two quotient
of $\pi_1(U_{\Q_p})$.

The ratio $C/x$ in (A.74.6) belongs to the uniquely
two-divisible group $1+p\Z_p$; hence
$\kappa_2(C)=\kappa_2(x)$.
The pro-two conjugacy therefore holds with the
cusp tangent $C$.
It maps to a conjugacy by $H_2$ under the actual
Legendre monodromy map, since $H_2$ is pro-two.
Together with the odd components and the product
description of $H$ in A.17, this proves full
$H$-conjugacy exactly as in (A.67.7).

It remains to determine the finite Kummer separation.
Equation (A.74.4), whose first nonconstant
coefficient is the $p$-adic unit $1/2$, gives
\[
                v_p(C/x-1)=e,
                     \qquad v_p((1-x)-1)=e.
\tag{A.74.8}
\]
The later terms have valuations at least $2e$.
For $k\ge0$ the exact power image is
\[
                (1+p\Z_p)^{p^k}=1+p^{k+1}\Z_p.
\tag{A.74.9}
\]
The root-of-unity part contributes no further
possibility: raising to $p^k$ is bijective on
$\mu_{p-1}$, so a unit of residue one can be
a $p^k$th power only through (A.74.9).
For $k<e$ both ratios comparing the abelian
coordinates $(x,1-x)$ and $(C,1)$ are therefore
$p^k$th powers.  At $k=e$ both are excluded.
This proves the assertion about the first
separating abelian $p$-power level.

The equation (A.74.2) defines a finite etale
$\mu_{p^e}$-torsor on $U$.
Its cusp-zero fibre has the rational sheet $z=1$,
so the cusp section lifts to the corresponding
open subgroup.  Its fibre at $x$ has no rational
sheet by (A.74.8)--(A.74.9); hence the point
section does not lift.
The geometric character has meridian values
$(0,1)$ at zero and one.
By A.56 it factors through $H$ exactly when
$p^e\mid24\cdot0$, $p^e\mid24\cdot1$ and
$p^e\mid8(0-1)$ simultaneously.
The last condition is impossible for odd $p$.
This proves the claimed noncongruence property,
including $p=3$ where $H^{\rm ab}$ does have a
diagonal cubic quotient.

Finally, in the abelian quotient modulo $p^e$,
cusp zero forces the second Kummer coordinate
to be trivial, contrary to (A.74.8).
Cusp one forces the first to be trivial, but
$v_p(x)=e$ is not divisible by $p^e$.
At infinity the two valuations must coincide,
whereas they are $e$ and zero and $p^e>e$.
Thus the pro-$p$ projection lies in no cusp packet.
For rational $x$, algebraicity of $C$ would make
$q_x=C^2/256$ algebraic, contradicting the proved
$p$-adic Mahler--Manin theorem because
$J(q_x)=j(E_x)$ is rational.  This proves the
last assertion as well.
\end{proof}

\noindent\emph{Exact finite certificate.}
The program \path{code/odd_local_congruence_certificate.py}
constructs the inverse nome series with rational arithmetic,
checks its coefficients and the square-root identity, and
tests the analytic isometry on finite odd-prime discs.
It also computes the Tate parameters for points of several
depths, checks both splitting characters, and checks the
exact first separating Kummer levels by modular arithmetic.
The retained output is
\path{code/odd_local_congruence_certificate_output.txt}.
These checks accompany the all-prime analytic and arithmetic
maps proved above.
